%% These are the bibliography items to be moved into the your bibtex file
\begin{filecontents*}{refs.bib}
@misc{vanbrunt_mcnew,
  author       = {Kai Van Brunt},
  title        = {Accuracy and Maneuverability of Atmospheric Reentry Vehicles},
  note         = {Unpublished report; modified MC-NEW aerodynamic workup for SWERVE},
  year         = {2024}
}

@techreport{langley1959,
  author       = {Armstrong, William O.},
  title        = {Static Longitudinal Aerodynamic Characteristics of Several Wing and Blunt-Body Shapes Applicable for Use as Reentry Configurations at a Mach Number of 6.8 and Angles of Attack up to 90 Deg},
  institution  = {National Aeronautics and Space Administration},
  number       = {NASA TM X-63},
  year         = {1959}
}

@techreport{maughmer1992control,
  author       = {Mark D. Maughmer and Lyle N. Long and Peter J. Pagano},
  title        = {Prediction of Forces and Moments for Hypersonic Flight Vehicle Control Effectors},
  institution  = {The Pennsylvania State University / NASA Langley Research Center},
  type         = {Semi-Annual Report},
  note         = {Grant No. NAG-849-1},
  year         = {1992}
}

@techreport{tauber1986,
  author       = {Michael E. Tauber},
  title        = {Maximum Lift/Drag Ratio of Flat Plates with Bluntness and Skin Friction at Hypersonic Speeds},
  institution  = {NASA Ames Research Center},
  number       = {NASA TM-88338},
  year         = {1986},
  month        = jul,
  note         = {NASA NTRS source available online}
}

@techreport{mayo1965,
  author       = {Edward E. Mayo},
  title        = {Newtonian Aerodynamics for Tangent Ogive Bodies of Revolution},
  institution  = {NASA Goddard Space Flight Center},
  number       = {X-671-65-244},
  year         = {1965},
  month        = jun
}

@misc{nasa_atmosphere,
  author       = {{NASA Glenn Research Center}},
  title        = {Earth Atmosphere Equation -- Metric},
  howpublished = {\url{https://www1.grc.nasa.gov/beginners-guide-to-aeronautics/earth-atmosphere-equation-metric/}},
  note         = {Accessed 2026-04-12},
  year         = {2026}
}

@misc{mason_hypersonics,
  author       = {William H. Mason},
  title        = {Configuration Aerodynamics for Hypersonic Vehicles},
  howpublished = {\url{https://archive.aoe.vt.edu/mason/Mason_f/ConfigAeroHypersonics.pdf}},
  note         = {Virginia Tech lecture notes; accessed 2026-04-12},
  year         = {n.d.}
}

@techreport{naca1135,
  author       = {{Ames Research Staff}},
  title        = {Equations, Tables, and Charts for Compressible Flow},
  institution  = {National Advisory Committee for Aeronautics (NACA)},
  number       = {Report 1135},
  year         = {1953}
}

@techreport{mcbride2002,
  author       = {McBride, Bonnie J. and Zehe, Michael J. and Gordon, Sanford},
  title        = {NASA Glenn Coefficients for Calculating Thermodynamic Properties of Individual Species},
  institution  = {NASA Glenn Research Center},
  number       = {NASA/TP-2002-211556},
  year         = {2002}
}

@techreport{gordon1994cea,
  author       = {Gordon, Sanford and McBride, Bonnie J.},
  title        = {Computer Program for Calculation of Complex Chemical Equilibrium Compositions and Applications. Part I: Analysis},
  institution  = {NASA Lewis Research Center},
  number       = {NASA RP-1311},
  year         = {1994}
}

@inproceedings{leader2024cea,
  author       = {Leader, Mark K. and Lavelle, Thomas M. and Wang, Xiao-yen J. and Dickens, Kevin W. and McTague, Michael and Hill, Jeffrey P.},
  title        = {CEA2022: A Modernization of NASA Glenn's Software CEA (Chemical Equilibrium with Applications)},
  booktitle    = {Thermal and Fluids Analysis Workshop Symposium},
  year         = {2024},
  note         = {NASA Glenn Research Center conference paper}
}

\end{filecontents*}

\documentclass{article}
\usepackage{graphicx} % Required for inserting images
\usepackage{amsmath}
\usepackage{xcolor}
\usepackage{url}
\usepackage{overarrows}

\title{Lift and Drag for Hypersonic Reentry}
\author{R. Scott Kemp}
\date{April 2026}

\begin{document}

\maketitle

\section{Purpose}
This is written for you, Amos. It is no longer a drop-in replacement for your write-up about how \texttt{pytrajlib} works, but we can strip out the implementation notes and use it for the publication (or documentation).

To start, \texttt{pytrajlib} should be modified so it swaps out RV parameters depending on whether maneuvering is turned on or not. When maneuvering is off, the RV should probably look something like a ballistic nuclear-weapon RV. For maneuvering flight (or ``terminal guidance''), the code should use SWERVE-like RV parameters. (In the future, the code could be made extensible for other non-SWERVE vehicles by allowing the user to specify which maneuverable reentry vehicle (MARV) is being used according to a pre-defined file format---provided the RV is axisymmetric and has four flaps so as not to require roll procedures (else, we are forced to do 6DOF commanded accelerations, which we rather avoid). In this version, all pitch, yaw, and roll effects are incidental and eventually stabilize away (although I suppose we might get some corkscrewing, not sure). I know that in some control codes there are explicit rules to ensure the angle-of-attack trends back to zero to avoid this. If, after implementing this, simulation accuracy degrades, plotting the angular rates to check for emergent corkscrewing should be an early diagnostic step.

All the vector and frame manipulations were developed with the help of AI and checked by another AI. I'm trusting they are correct. 

\section{The aerodynamic parameters}
Aerodynamic parameters are the nondimensional coefficients (and their derivatives)---\(C_D\), \(C_L\), \(C_N\), \(C_M\), and \(C_{M_q}\)---that quantify how a vehicle's shape and attitude convert airflow into forces and moments (torques). These inputs are used by \texttt{pytrajlib} to compute trajectory and perform flight control, and they control the vehicle's trim and stability. In practice, all of these parameters are functions of the vehicle's angle of attack (\(\alpha\)) relative to the freestream vector of oncoming air, which is formally defined in section~\ref{sec:baselineaero}.

Using a reference area \(S_{\mathrm{ref}}\) (typically in \(\mathrm{m}^2\)) and reference length \(c_{\mathrm{ref}}\) (typically in \(\mathrm{m}\)), the dimensional aerodynamic forces and moments acting on the body are written as
\begin{align}
D &= q_\infty S_{\mathrm{ref}} C_D, \label{eq:drag_force} \\
L &= q_\infty S_{\mathrm{ref}} C_L, \label{eq:lift_force} \\
N &= q_\infty S_{\mathrm{ref}} C_N, \label{eq:normal_force} \\
M &= q_\infty S_{\mathrm{ref}} c_{\mathrm{ref}} C_M, \label{eq:pitch_moment}
\end{align}
where 
\begin{equation}
q_\infty = \frac{1}{2}\rho V^2 \label{eq:dynamic_pressure}
\end{equation}
is the freestream dynamic pressure of the oncoming air before it begins to interact with the vehicle, \(\rho\) is atmospheric density at the current altitude and \(V\) is the vehicle airspeed. (\textbf{Implementation Note:} ensure strict consistency of metric units; assuming MKS (meter-kilogram-second) units, the forces should be evaluated in Newtons and the moments in \(\mathrm{N\cdot m}\).)

The coefficients are as follows: 
\begin{itemize}
    \item \(C_D\) is the drag coefficient, representing the aerodynamic force component opposite the freestream direction. It determines the axial deceleration of the vehicle and the drag penalty associated with flap actuation.
    \item \(C_L\) is the lift coefficient, representing the force component perpendicular to the freestream in the maneuver plane; it sets the lateral acceleration available for guidance and control.
    \item \(C_N\) is the normal-force coefficient, defined with respect to the body axis rather than the wind axis. For the nearly axisymmetric conical configurations considered here, \(C_N\) is the coefficient most directly associated with angle-of-attack loading and is often close in magnitude to \(C_L\) at small angles of attack.
    \item \(C_M\) describes the aerodynamic pitching moment about the center of gravity. It is used to determine trim and static stability. In particular, the trim condition is \(C_M = 0\), and the restoring (or destabilizing) tendency of the vehicle is measured by \(C_{M_\alpha} = \frac{\partial C_M}{\partial \alpha}\).
    \item \(C_{M_q}\) is the pitch-damping derivative. It gives the additional aerodynamic moment produced by a pitch rate \(q\), namely \(M_{\mathrm{damp}} = \frac{1}{4}\rho V S_{\mathrm{ref}} c_{\mathrm{ref}}^2 C_{M_q}\, q\). A negative value of \(C_{M_q}\) corresponds to aerodynamic damping and causes angular disturbances to decay in time.
\end{itemize}  

The aerodynamic coefficients are used in two ways. First, \(C_D\), \(C_L\), and \(C_N\) determine the translational aerodynamic forces appearing in the trajectory and guidance model. Second, \(C_M\) and \(C_{M_q}\) determine the rotational aerodynamics appearing in the attitude-dynamics model. In the linearized pitch equation, these effects enter through
\begin{equation}
\ddot{\alpha}
-
\frac{\rho V S_{\mathrm{ref}} c_{\mathrm{ref}}^2}{4I_y} C_{M_q}\dot{\alpha}
-
\frac{\rho V^2 S_{\mathrm{ref}} c_{\mathrm{ref}}}{2I_y} C_{M_\alpha}\alpha
= 0, \label{eq:linear_pitch}
\end{equation}
where \(I_y\) is the pitch-axis moment of inertia (in \(\mathrm{kg\cdot m^2}\)).


\section{Computation of Aerodynamic parameters for a SWERVE-like vehicle}
The Sandia Winged Energetic Re-entry Vehicle Experiment (SWERVE) was a maneuverable hypersonic vehicle developed by the United States between 1970 and 1985, which we adopt as our template maneuverable reentry vehicle. Kai Van Brunt performed a practical hypersonic aerodynamic workup for a conical reentry vehicle modeled on SWERVE by modifying the open-source code \texttt{MC-NEW}. This code computes aerodynamic coefficients through Monte Carlo integration of Newtonian-impact pressure over the vehicle surface, using a body reference area and length. The method assumes a simplified geometric representation of the vehicle and flaps, models local pressure with Newtonian theory plus a shock-correction factor for supersonic flow, approximates flap cross-sections and normals with thin-surface idealizations, and estimates pitch damping from a small-angular-rate perturbation to the local flow. The model explicitly neglects several effects, including turbulence, ablation, and high-temperature real-gas phenomena. However, in the hypersonic regime, the model is a reasonable approximation of what might be encountered with a SWERVE-like vehicle: Van Brunt shows that it reproduces analytical cone results for \(C_N\), \(C_M\), and \(C_{M_q}\) to within a few percent in regimes where the theory applies, but agreement with higher-fidelity Parabolized Navier–Stokes results for the full SWERVE vehicle is only moderate, with errors on the order of 20--30\%. Since our purpose is to approximate the behavior, rather than simulate an actual vehicle, we judge this to be sufficient---but this work is not as a substitute for flight-calibrated computational fluid dynamics that would be required for a real-world design. 

For the SWERVE geometry used by Van Brunt, the reference area for the undeflected-flap vehicle is taken as the cross-sectional base area of the conical body, and the reference length is taken as the total tip-to-base vehicle length:
\begin{equation}
S_{\mathrm{ref}}=\pi R_b^2,
\qquad
c_{\mathrm{ref}}=L, \label{eq:reference_dims}
\end{equation}
where \(R_b\) is the cone base radius and \(L\) is the overall vehicle length. 
\[
R_b = 0.277~\mathrm{m},
\qquad
L = 2.75~\mathrm{m},
\]
so that
\[
S_{\mathrm{ref}} \approx \pi(0.277)^2 \approx 0.241~\mathrm{m}^2,
\qquad
c_{\mathrm{ref}} = 2.75~\mathrm{m}.
\]

Van Brunt performed these calculations at Mach 12 across a range of angles of attack (\(\alpha\)), but with all flaps in their undeflected state. For the present application, the parameters computed at Mach-12 can be used over the interval of interest \(8 \le M_\infty \le 25\). The justification for extending the range of validity is that, in modified Newtonian theory, pressure-derived coefficients scale with the stagnation-point pressure coefficient \(C_{p,\max}(M_\infty)\), so the Mach dependence enters chiefly through
\begin{equation}
r(M_\infty) \equiv \frac{C_{p,\max}(M_\infty)}{C_{p,\max}(12)},
\qquad
C_{p,\max}(M_\infty)=\frac{2}{\gamma M_\infty^2}\left(\frac{p_{0,2}}{p_\infty}-1\right), \label{eq:mach_scaling}
\end{equation}
where \(r\) is a scaling factor for different Mach numbers of the freestream air, \(p_{0,2}\) is the post-shock stagnation pressure from the normal-shock relations~\cite{naca1135,langley1959}, \(p_\infty\) is the freestream static pressure, and \(\gamma\) is the heat-capacity ratio of air. 

For a calorically perfect gas, \(\frac{p_{0,2}}{p_\infty}\) can be written
\begin{equation}
\frac{p_{0,2}}{p_\infty}
=
\left(1+\frac{\gamma-1}{2}M_\infty^2\right)^{\frac{\gamma}{\gamma-1}}
\left[
\frac{\gamma+1}{2\gamma M_\infty^2-(\gamma-1)}
\right]^{\frac{1}{\gamma-1}}
\left[
\frac{(\gamma+1)M_\infty^2}{(\gamma-1)M_\infty^2+2}
\right]^{\frac{\gamma}{\gamma-1}}, \label{eq:post_shock_pressure}
\end{equation}
and we can set \(\gamma=1.4\). 

Evaluating the above relations over the Mach-number range of interest, we find that the resulting error in the stagnation-point pressure coefficient \(C_{p,\max}\) is estimated to be small. In particular, varying the Mach or \(\gamma\) over a representative hypersonic range \(8 \le M_\infty \le 22\) and \(1.15 \le \gamma \le 1.40\) changes the perfect-gas value of \(C_{p,\max}\) by less than \(5\%\). This bound does not, however, include high-temperature real-gas effects such as vibrational excitation, dissociation, or chemical nonequilibrium. Those effects can become important in faster hypersonic flows, especially above about Mach 15. In Appendix~\ref{app:realgas}, we additionally show that for equilibrium hot real gas, the effects are still small and the calorically perfect, constant-Mach assumption, \(r(M_\infty)\approx1\) will hold over the range of interest. For interest, the computed equilibrium real-gas values for \(r(M_\infty)\) are tabulated in Table~\ref{tab:realgascorrection} and can be incorporated into the code with linear interpolation if you desire.

\subsection{Baseline body aerodynamics, reduced state, and derived quantities}\label{sec:baselineaero}

The original state \((\hat{\mathbf b},\alpha,\dot{\alpha})\) was adequate only for a single longitudinal control axis. Once two orthogonal flap pairs and arbitrary azimuth are introduced, a scalar angle of attack no longer determines (i) the direction of the wind in the body-transverse plane, (ii) the azimuth of the lift vector, or (iii) which flap pair is most exposed to the freestream. We therefore have to expand the model by adding a full body-centric frame that can be reconstructed at every time step. Although 6DOF can be resolved in flight, this is still not capable of supporting a full 6-DOF navigation model as the body is treated as axisymmetric with four flaps. Roll commands can be processed by this procedure, but at present there are no plans to support them in the Autopilot function, so we anticipate \(\omega_{3,B}\) will always be zero.

Let
\begin{equation}
\mathcal E=\{\hat{\mathbf e}_{x,E},\hat{\mathbf e}_{y,E},\hat{\mathbf e}_{z,E}\}
\end{equation}
denote the Earth-centered inertial (ECI) basis, and let
\begin{equation}
\mathcal B=\{\hat{\mathbf e}_{1,B},\hat{\mathbf e}_{2,B},\hat{\mathbf e}_{3,B}\},
\qquad
\hat{\mathbf e}_{3,B}=
\begin{bmatrix}
0\\[2pt]0\\[2pt]1
\end{bmatrix},
\end{equation}
denote the body-fixed basis, with \(\hat{\mathbf e}_{3,B}\) aligned with the cone axis and directed towards the base of the vehicle, so that \(+\hat{\mathbf e}_{3,B}\) points aft. With this convention, \(\alpha=0\) when the relative-wind direction \(\hat{\mathbf u}_B\) is equal to \(+\hat{\mathbf e}_{3,B}\), i.e. when the airflow in body-centric coordinates is directed from the tip toward the base. (\textbf{Implementation Note:} To complete a right-handed system, \(+\hat{\mathbf e}_{1,B}\) and \(+\hat{\mathbf e}_{2,B}\) typically form the transverse axes, e.g., ``dorsal'' and ``starboard'' respectively.) 

The body axis in ECI coordinates is recovered as
\begin{equation}
\hat{\mathbf b}_E = C_{EB}\hat{\mathbf e}_{3,B}, \label{eq:body_axis_eci}
\end{equation}
where \(C_{EB}\in SO(3)\) is a \(3\times3\) Direction Cosine Matrix that maps body-frame components into ECI components, and \(C_{BE}=C_{EB}^{\mathsf T}\) maps ECI components into body-frame components. A key point here is that \(\mathcal B\) moves with the vehicle and keeps track of the vehicle's attitude, therefore \(C_{EB}\) and its transpose \(C_{BE}\) change with every time step and must be continuously recomputed. 

In simpler 3DOF or restricted 6DOF models, one might attempt to track attitude using just the symmetry axis $\hat{\mathbf b}_E$, a transverse reference vector, and an explicit scalar roll angle \(\phi\) to avoid gimbal lock. However, maintaining strict orthonormality constraints ($C_{EB}^{\mathsf T}C_{EB}=I, \det C_{EB} =1$) through numerical integration over many time steps is expensive and prone to error. By contrast, a unit quaternion avoids singularities intrinsically and only requires a simple length constraint:
\begin{equation}
\|\mathbf q\| = 1.
\end{equation}
This constraint should be enforced numerically at the end of each step via simple renormalization:
\begin{equation}
\mathbf q \leftarrow \frac{\mathbf q}{\|\mathbf q\|}. \label{eq:quat_renorm}
\end{equation}

Thus, we define the single-unit quaternion  
\begin{equation}
\mathbf q_{EB}=
\begin{bmatrix}
q_0\\ q_1\\ q_2\\ q_3
\end{bmatrix},
\qquad
\|\mathbf q_{EB}\|=1, \label{eq:quaternion_def}
\end{equation}
where \(q_0\) is the cosine of the half-angle of rotation, and \((q_1, q_2, q_3)\) define a vector that points along the axis of rotation and has magnitude that is the sine of the half angle of rotation. The full quaternion \(\mathbf q_{EB}\) can be interpreted so that it natively generates the full 3D body-to-ECI mapping. (\textbf{Implementation Note:} Ensure any third-party spatial library uses the same convention assumed here: scalar-first quaternion storage, Hamilton multiplication, and an active body-to-ECI mapping---all three are required. Libraries that use scalar-last storage, passive rotations, or the opposite multiplication order require corresponding adjustments.)

\subsubsection{State Vector}
Using the quaternion, the required state vector for the vehicle is then
\begin{equation}
\mathbf x_q =
\begin{bmatrix}
\mathbf r_E\\
\mathbf v_E\\
\mathbf q_{EB}\\
\boldsymbol\omega_{B}\\
\delta_1\\
\delta_2
\end{bmatrix}. \label{eq:state_vector}
\end{equation}
Here, the vector
\begin{equation}
\boldsymbol\omega_B=
\begin{bmatrix}
\omega_{1,B}\\[2pt]\omega_{2,B}\\[2pt]\omega_{3,B}
\end{bmatrix}
\end{equation}
is the vehicle angular velocity written in body components: \(\omega_{1,B}\) is the instantaneous rotation rate about the body axis \(\hat{\mathbf e}_{1,B}\), \(\omega_{2,B}\) about \(\hat{\mathbf e}_{2,B}\), and \(\omega_{3,B}\) about the symmetry axis \(\hat{\mathbf e}_{3,B}\). 

The quantities \(\delta_1,\delta_2\) are the two flap-pair deflections. The commanded flap deflections are defined by the right-hand rule about the corresponding hinge axes. Let flap pair 1 consist of flaps \(\{1,3\}\) with common hinge axis
\begin{equation}
\hat{\mathbf h}_{1,B}=\hat{\mathbf h}_{3,B}=\hat{\mathbf e}_{2,B},
\end{equation}
and let flap pair 2 consist of flaps \(\{2,4\}\) with common hinge axis
\begin{equation}
\hat{\mathbf h}_{2,B}=\hat{\mathbf h}_{4,B}=-\hat{\mathbf e}_{1,B}.
\end{equation}
The sign convention for flap deflection follows the right-hand rule about the hinge axis \(\hat{\mathbf h}_{i,B}\): if the right-hand thumb points along \(\hat{\mathbf h}_{i,B}\), then the curl of the fingers gives the positive direction of flap rotation. Therefore \(+\delta_1\) denotes a positive rotation of flaps \(1\) and \(3\) about \(+\hat{\mathbf e}_{2,B}\), and \(+\delta_2\) denotes a positive rotation of flaps \(2\) and \(4\) about \(-\hat{\mathbf e}_{1,B}\). With this convention, the sign of the aerodynamic increment produced by a positive \(\delta_p\) is determined by the geometry through the rotated flap normals and the resulting flap-force calculation, rather than by a separate sign inserted by hand.

In this implementation, the primary body axis \(\hat{\mathbf b}_E\) and the full direction cosine matrix \(C_{EB}\) are directly reconstructed from \(\mathbf q_{EB}\) at each time step as detailed in Section~\ref{sec:frametrans}.

\subsubsection{Tracking Roll}
It is worth noting that there are two dynamic effects often colloquially grouped as "roll." One phenomenon is an explicit twisting rotation of the body about its symmetry axis \(\hat{\mathbf e}_{3,B}\), represented by \(\omega_{3,B}\). We do not currently plan to command such body-axial rolls.

A second phenomenon often informally called ``roll'' is the geometric bank or rotation of the body-transverse axes that arises when the symmetry axis changes direction through noncommutative three-dimensional rotations such as pitch$\longrightarrow$yaw$\longrightarrow$pitch. 

The quaternion attitude state \(\mathbf q_{EB}\) and its derived quantities $C_{EB}$ and $C_{BE}$ track both the explicit axial twist and this geometric banking. In the present reduced model, however, body-axial twist is intentionally suppressed. Accordingly, Step~9 enforces
\[
\omega_{3,B}=0,
\qquad
M_{3,B}=0,
\]
so that only the geometric banking evolves dynamically. If explicit axial-roll dynamics are introduced in the future, Step~9 must use the full axisymmetric Euler equations.

\subsection{Winds in the Body Frame}
The relative wind is
\begin{equation}
\mathbf V_{\mathrm{rel},E} = \mathbf V_{\mathrm{wind},E} - \mathbf v_E,
\qquad
V = \|\mathbf V_{\mathrm{rel},E}\|,
\qquad
\hat{\mathbf u}_E = \frac{\mathbf V_{\mathrm{rel},E}}{V}, \label{eq:relative_wind_eci}
\end{equation}
and in body components
\begin{equation}
\mathbf V_{\mathrm{rel},B}=C_{BE}\mathbf V_{\mathrm{rel},E},
\qquad
\hat{\mathbf u}_B=C_{BE}\hat{\mathbf u}_E=
\begin{bmatrix}
u_1\\u_2\\u_3
\end{bmatrix}. \label{eq:relative_wind_body}
\end{equation}
The local dynamic pressure and Mach number are
\begin{equation}
q_\infty=\tfrac12 \rho(\mathbf r_E)V^2,
\qquad
M_\infty=\frac{V}{a_\infty(\mathbf r_E)}. \label{eq:mach_def}
\end{equation}
\textbf{Implementation Note:} Ensure everywhere the Mach number is calculated in the code, it is dynamically referenced to the local speed of sound \(a_\infty\). This can come directly from \texttt{EarthGram} or you can compute it via \(a_\infty = \sqrt{\gamma_\infty R T_\infty}\), where \(\gamma_\infty = 1.4\) and \(T_\infty\) is the local, unperturbed air temperature. Additionally, be careful of vector normalizations. While the velocity normalization above is likely always going to be fine, in general if the norm of the vector being normalized is below a small tolerance, the associated direction is undefined. In that event, treat the state as outside the intended operating regime: skip the directional calculation for that step and set the corresponding aerodynamic force and moment contributions to zero. Sine we are almost never going to be in the low-speed regime, this shouldn't be an issue, however.

Every externally supplied vector is used with its literal physical sign. For example, \(\mathbf g_E\) is the actual gravitational acceleration vector, so at a point on the positive \(\hat{\mathbf e}_{z,E}\) axis one has approximately
\[
\mathbf g_E\approx -g\,\hat{\mathbf e}_{z,E}.
\]
Likewise \(\mathbf V_{\mathrm{wind},E}\) is the velocity of the air mass itself, not the direction ``from which'' the wind comes. Thus if the air is locally still in ECI coordinates, \(\mathbf V_{\mathrm{wind},E}=0\) and
\[
\mathbf V_{\mathrm{rel},E}=-\mathbf v_E,
\qquad
\hat{\mathbf u}_B=\hat{\mathbf e}_{3,B}
\quad\text{when }\alpha=0.
\]
With this convention, \(\hat{\mathbf u}\) points in the direction of the relative wind (i.e., aft towards the base), so a positive drag coefficient multiplies \(+\hat{\mathbf u}_B\), which generates force opposite the vehicle motion in still air.

The actual angle of attack is most cleanly computed from \(\hat{\mathbf u}_B\):
\begin{equation}
\alpha=\cos^{-1}(u_3), \label{eq:aoa_scalar}
\end{equation}
taking care to address numerical round-off errors by clamping 
\(u_3\leftarrow\textrm{clip}(u_3,-1,1)\).
The azimuth of the attack vector in the body-transverse plane is
\begin{equation}
\chi_\alpha=\operatorname{atan2}(u_2,u_1), \label{eq:aoa_azimuth}
\end{equation}
taking care to note the azimuth \(\chi_\alpha\) is undefined when \(u_1=u_2=0\),
and a convenient vector measure of the instantaneous attack state is
\begin{equation}
\boldsymbol\alpha_{\perp,B}
=
\frac{1}{u_3}
\begin{bmatrix}
u_1\\[2pt]u_2
\end{bmatrix}
=
\tan\alpha
\begin{bmatrix}
\cos\chi_\alpha\\[2pt]\sin\chi_\alpha
\end{bmatrix}, \label{eq:aoa_vector}
\end{equation}
which, for small angles, reduces componentwise to the usual pitch- and yaw-like attack components. This vector is the correct replacement for a single scalar \(\alpha\) in the two-pair problem. 

\textbf{Implementation Note:} \(u_3 \approx 0\) would cause a singularity, but a conical reentry vehicle will be aerodynamically unstable long before \(\alpha\) approaches \(90^\circ\), making this division physically safe. However, in code, this is a good place for a numeric assertion to safely catch arbitrary tumbling states. To avoid the azimuth \(\chi_\alpha\) from becoming undefined (\(u_1=u_2=0\)), test for \( s=\sqrt{u_1^2+u_2^2}<\epsilon_s,\), and if true set \(\chi_\alpha\) to zero or retain its previous value only as a placeholder for continuity. 

Define the body-lift direction in body coordinates by projecting the body axis into the plane normal to the freestream:
\begin{equation}
\hat{\boldsymbol\ell}_B
=
\frac{\left(\mathbf I-\hat{\mathbf u}_B\hat{\mathbf u}_B^{\mathsf T}\right)\hat{\mathbf e}_{3,B}}
{\left\|\left(\mathbf I-\hat{\mathbf u}_B\hat{\mathbf u}_B^{\mathsf T}\right)\hat{\mathbf e}_{3,B}\right\|}, \label{eq:lift_direction}
\end{equation}
(\textbf{Implementation Note:} The vector \(\left(\mathbf I-\hat{\mathbf u}_B\hat{\mathbf u}_B^{\mathsf T}\right)\hat{\mathbf e}_{3,B}\) takes the aft-pointing body axis and subtracts its projection along the relative wind. Since the body is physically tilted \emph{away} from the windward side, this resultant vector rigorously points toward the \emph{leeward} side. Thus, it correctly gives the physical direction of positive lift away from the windward face.)

Define the body pitching-moment axis by
\begin{equation}
\hat{\mathbf m}_B
=
\frac{\hat{\mathbf u}_B\times \hat{\mathbf e}_{3,B}}
{\|\hat{\mathbf u}_B\times \hat{\mathbf e}_{3,B}\|}. \label{eq:moment_axis}
\end{equation}
The matrix \(\mathbf I-\hat{\mathbf u}_B\hat{\mathbf u}_B^{\mathsf T}\) is the projector onto the plane perpendicular to the wind. It removes from \(\hat{\mathbf e}_{3,B}\) the part that lies parallel to the wind and leaves only the sideways component. After normalization, \(\hat{\boldsymbol\ell}_B\) therefore points in the body-fixed direction of positive lift. The cross product then chooses the corresponding pitching-moment axis. As a sign check, if the attack vector lies in the \((\hat{\mathbf e}_{1,B},\hat{\mathbf e}_{3,B})\) plane with \(u_1>0\) and \(u_2=0\), then \(\hat{\mathbf m}_B=-\hat{\mathbf e}_{2,B}\); in that case a positive tabulated \(C_M\) produces a moment in the \(-\hat{\mathbf e}_{2,B}\) direction.

When \(\alpha\) is numerically very small, both vector directions in equations \eqref{eq:lift_direction} and \eqref{eq:moment_axis} become undefined because their denominators scale as \(\sin\alpha\). These directions only multiply the baseline quantities that vanish at zero angle of attack, namely \(C_L(0)=0\) and \(C_M(0)=0\). In code, you can detect this state by checking for the condition 
\begin{equation}
s=\sqrt{u_1^2+u_2^2}<\epsilon.    
\end{equation}
When this occurs, set the baseline (non-flap) lift and static-moment contributions to be zero. If this condition fails (angle-of-attack is nontrivial), then compute: 
\begin{equation}
\hat{\mathbf m}_B = \frac{1}{s}
\begin{bmatrix}
u_2 \\
-\,u_1 \\
0
\end{bmatrix},
\qquad
\hat{\boldsymbol \ell}_B = \frac{1}{s}
\begin{bmatrix}
-\,u_1 u_3 \\
-\,u_2 u_3 \\
1 - u_3^2
\end{bmatrix}.
\end{equation}

Their ECI images can be recovered:
\begin{equation}
\hat{\boldsymbol\ell}_E=C_{EB}\hat{\boldsymbol\ell}_B,
\qquad
\hat{\mathbf m}_E=C_{EB}\hat{\mathbf m}_B.
\end{equation}

With the undeflected-vehicle aerodynamic tables, the baseline body force in body coordinates is
\begin{equation}
\mathbf F_{\mathrm{body},B}
=
q_\infty S_{\mathrm{ref}}
\left[
C_D(\alpha,M_\infty)\hat{\mathbf u}_B
+
C_L(\alpha,M_\infty)\hat{\boldsymbol\ell}_B
\right], \label{eq:body_force_baseline}
\end{equation}
\begin{equation}
\mathbf M_{\mathrm{body},B}^{(\mathrm{stat})}
=
q_\infty S_{\mathrm{ref}} c_{\mathrm{ref}}
\,C_M(\alpha,M_\infty)\hat{\mathbf m}_B. \label{eq:body_moment_baseline}
\end{equation}
The vehicle angular velocity is already stored in body components. Its transverse part,
\begin{equation}
\boldsymbol\omega_{\perp,B}
=
\begin{bmatrix}
\omega_{1,B}\\[2pt]\omega_{2,B}\\[2pt]0
\end{bmatrix},
\end{equation}
contains the rates that tip the body axis away from the freestream and therefore change the angle of attack. The aerodynamic damping contribution is written in terms of this transverse part so that
\begin{equation}
\mathbf M_{\mathrm{body},B}^{(\mathrm{damp})}
=
q_\infty S_{\mathrm{ref}} c_{\mathrm{ref}}
\left(\frac{c_{\mathrm{ref}}}{2V}\right)
C_{M_q}(\alpha,M_\infty)\,\boldsymbol\omega_{\perp,B}. \label{eq:aero_damping}
\end{equation}
Thus, the total baseline body moment is
\begin{equation}
\mathbf M_{\mathrm{body},B}
=
\mathbf M_{\mathrm{body},B}^{(\mathrm{stat})}
+
\mathbf M_{\mathrm{body},B}^{(\mathrm{damp})}.
\end{equation}
The signs are carried by the tabulated coefficients themselves; for a statically stable, damped vehicle one typically has \(C_M<0\) for positive \(\alpha\) and \(C_{M_q}<0\).

The translational equations of motion are
\begin{equation}
\dot{\mathbf r}_E=\mathbf v_E,
\qquad
m\dot{\mathbf v}_E
=
m\mathbf g_E(\mathbf r_E)+\mathbf F_{\mathrm{aero},E}, \label{eq:eom_translation}
\end{equation}
with
\begin{equation}
\mathbf F_{\mathrm{aero},E}=C_{EB}\mathbf F_{\mathrm{aero},B}.
\end{equation}
The reduced rotational equations are
\begin{equation}
I_\perp \dot{\omega}_{1,B}=M_{1,B},
\qquad
I_\perp \dot{\omega}_{2,B}=M_{2,B},
\qquad
I_3 \dot{\omega}_{3,B}=M_{3,B}, \label{eq:eom_rotation}
\end{equation}
where \(I_\perp\) is the transverse moment of inertia and \(I_3\) is the axial moment of inertia. 

Your original one-axis linearized rotational equation written in terms of \(\alpha\) is accurate about a fixed maneuver plane:
\[
I_\perp \ddot{\alpha}
\approx
q_\infty S_{\mathrm{ref}} c_{\mathrm{ref}}
\left(
C_{M_\delta}\delta
+
C_{M_\alpha}\alpha
+
C_{M_q}\frac{c_{\mathrm{ref}}}{2V}\dot{\alpha}
\right).
\]
However, for the new two-flap-pair problem, the correct linearized generalization is the vector equation
\begin{equation}
I_\perp \ddot{\boldsymbol\alpha}_{\perp,B}
\approx
q_\infty S_{\mathrm{ref}} c_{\mathrm{ref}}
\left(
C_{M_\alpha}\boldsymbol\alpha_{\perp,B}
+
C_{M_q}\frac{c_{\mathrm{ref}}}{2V}\dot{\boldsymbol\alpha}_{\perp,B}
+
C_{M_\delta}\boldsymbol\delta_{\mathrm{eff}}
\right), \label{eq:linear_pitch_vector}
\end{equation}
with
\begin{equation}
\boldsymbol\delta_{\mathrm{eff}}
=
\begin{bmatrix}
\delta_{1,\mathrm{eff}}\\[2pt]\delta_{2,\mathrm{eff}}
\end{bmatrix}.
\end{equation}
Writing
\[
p=\frac{C_{M_q}q_\infty S_{\mathrm{ref}}c_{\mathrm{ref}}^2}{2I_\perp V},
\qquad
k=\frac{C_{M_\alpha}q_\infty S_{\mathrm{ref}}c_{\mathrm{ref}}}{I_\perp},
\qquad
n=\frac{C_{M_\delta}q_\infty S_{\mathrm{ref}}c_{\mathrm{ref}}}{I_\perp},
\]
one obtains
\begin{equation}
\dot{\boldsymbol\alpha}_{\perp,B}=\boldsymbol\nu_B,
\qquad
\dot{\boldsymbol\nu}_B
=
p\,\boldsymbol\nu_B
+
k\,\boldsymbol\alpha_{\perp,B}
+
n\,\boldsymbol\delta_{\mathrm{eff}}.
\end{equation}
This form is useful for controller design, but the simulation loop below uses the full force and moment vectors because those remain valid away from the linearized trim point.

\subsection{Frame Translation and the Direction Cosine Matrix}\label{sec:frametrans}

The Direction Cosine Matrix \(C_{EB}\) is the primary bridge between the rigid body and the simulated universe. It converts vectors in the body frame \(\mathcal{B}\) to the ECI frame \(\mathcal{E}\), and its transpose \(C_{BE} = C_{EB}^{\mathsf T}\) does the reverse:
\begin{equation}
\mathbf v_E=C_{EB}\mathbf v_B,
\qquad
\mathbf v_B=C_{BE}\mathbf v_E.
\end{equation}

The matrix \(C_{EB}\) is constructed such that its three columns are the body-frame basis vectors expressed in ECI coordinates:
\begin{equation}
C_{EB}
=
\begin{bmatrix}
[\hat{\mathbf e}_{1,B}]_E &
[\hat{\mathbf e}_{2,B}]_E &
[\hat{\mathbf e}_{3,B}]_E
\end{bmatrix},
\qquad
[\hat{\mathbf e}_{3,B}]_E=\hat{\mathbf b}_E.
\label{eq:dcm_definition}
\end{equation}

Because propagating \(C_{EB}\) directly via rotational kinematics causes the columns to slowly lose their orthonormality, we only maintain the state of the system via the unit quaternion $\mathbf{q}_{EB}$. At any step where vector transformations or aerodynamic measurements are needed, \(C_{EB}\) is cleanly mapped back from the quaternion. 

Given our scalar-first unit quaternion state vector $\mathbf{q}_{EB} = [q_0, q_1, q_2, q_3]^{\mathsf T}$, the exact, computationally safe equivalent transformation matrix is:
\begin{equation}
C_{EB} = \begin{bmatrix}
1 - 2(q_2^2 + q_3^2) & 2(q_1q_2 - q_0q_3) & 2(q_1q_3 + q_0q_2) \\[4pt]
2(q_1q_2 + q_0q_3) & 1 - 2(q_1^2 + q_3^2) & 2(q_2q_3 - q_0q_1) \\[4pt]
2(q_1q_3 - q_0q_2) & 2(q_2q_3 + q_0q_1) & 1 - 2(q_1^2 + q_2^2)
\end{bmatrix}. \label{eq:quat_to_dcm}
\end{equation}
recalling that
\begin{equation}
C_{BE}=C_{EB}^{\mathsf T}    
\end{equation}

By keeping this mapping strictly localized, the simulation enjoys the geometric robustness of quaternions without losing the straightforward math of linear algebra mapping.

The intuitive body-centric parameters can be recovered as follows. The ECI representation of the cone axis is simply
\[
[\hat{\mathbf e}_{3,B}]_E = \hat{\mathbf b}_E = C_{EB}\hat{\mathbf e}_{3,B}.
\]

No separate scalar roll angle is required by the allocator or aerodynamic model: any body-axis twist and any geometric banking are already contained in \(C_{EB}\) (equivalently in \(\mathbf q_{EB}\)). If a scalar body-axis twist angle is ever desired for diagnostics, it must be defined relative to a chosen inertial reference vector \(\mathbf d_E\) that is not parallel to \(\hat{\mathbf b}_E\). One convenient construction is
\[
\hat{\mathbf a}_E =
\frac{(\mathbf I - \hat{\mathbf b}_E \hat{\mathbf b}_E^{T}) \, \mathbf d_E}
{\left\| (\mathbf I - \hat{\mathbf b}_E \hat{\mathbf b}_E^{T}) \, \mathbf d_E \right\|},
\qquad
\hat{\mathbf c}_E = \hat{\mathbf b}_E \times \hat{\mathbf a}_E,
\]
and with
\[
[\hat{\mathbf e}_{1,B}]_E = C_{EB}\hat{\mathbf e}_{1,B}
\]
one may define the diagnostic twist angle
\[
\phi = \operatorname{atan2}
\!\left(
[\hat{\mathbf e}_{1,B}]_E \cdot \hat{\mathbf c}_E,\;
[\hat{\mathbf e}_{1,B}]_E \cdot \hat{\mathbf a}_E
\right).
\]
This \(\phi\) is only a diagnostic angle relative to the chosen reference \(\mathbf d_E\); it is not needed for control allocation.

%Because the body is fourfold symmetric, the control allocation is periodic in \(\phi\) with period \(\pi/2\) and can be reduced modulo \(\pi/2\) inside the allocator if desired.

\section{Workable maneuver loop for two orthogonal flap pairs}

The loop below follows our proposed sequence. The key conceptual change is that the attitude update is performed with the vehicle angular-velocity vector expressed in body coordinates, after which \(\alpha\) is recomputed as a derived quantity.

\subsection{Step 1: Preprocess the commanded acceleration}

Let the guidance law produce an inertial acceleration command \(\mathbf a_{\mathrm{cmd},E}\) (in \(\mathrm{m/s^2}\)) that is intended to be transverse to the estimated direction of motion, but which does not account for wind. Define the estimated velocity direction seen by guidance as
\begin{equation}
\hat{\mathbf v}_{\mathrm{est},E}
=
\frac{\mathbf v_E}{\|\mathbf v_E\|}.
\end{equation}
If the guidance command includes gravity, remove it first:
\begin{equation}
\mathbf a_{\mathrm{aero,cmd},E}^{\star}
=
\mathbf a_{\mathrm{cmd},E}-\mathbf g_E(\mathbf r_E). \label{eq:commanded_accel_star}
\end{equation}
(\textbf{Implementation Note:} The sign convention depends on whether the guidance law outputs total required kinematic acceleration or aerodynamic acceleration. Ensure literal vector addition is used here, respecting the actual signs of the ECI vectors rather than subtracting magnitudes. Here again we also have a normalization that can go badly if the velocity becomes too low.) 
For example, if \(\mathbf a_{\mathrm{cmd},E}=2\,\hat{\mathbf e}_{x,E}\,\mathrm{m/s^2}\) and the local gravitational acceleration is \(\mathbf g_E=-9.6\,\hat{\mathbf e}_{z,E}\,\mathrm{m/s^2}\), then
\[
\mathbf a_{\mathrm{aero,cmd},E}^{\star}=2\,\hat{\mathbf e}_{x,E}+9.6\,\hat{\mathbf e}_{z,E}\,\mathrm{m/s^2}
\]
before the transverse projection is applied. Here ``transverse'' simply means perpendicular to the estimated direction of motion. Project this command into that transverse plane:
\begin{equation}
\mathbf a_{\mathrm{aero,cmd},E}
=
\left(
\mathbf I-\hat{\mathbf v}_{\mathrm{est},E}\hat{\mathbf v}_{\mathrm{est},E}^{\mathsf T}
\right)
\mathbf a_{\mathrm{aero,cmd},E}^{\star}. \label{eq:commanded_accel_projected}
\end{equation}
If the guidance law already outputs non-gravitational transverse acceleration, then this step reduces to the identity.

\subsection{Step 2: Autopilot allocation to the two flap pairs}

Let
\begin{equation}
P_{\perp,E}
=
\mathbf I-\hat{\mathbf v}_{\mathrm{est},E}\hat{\mathbf v}_{\mathrm{est},E}^{\mathsf T}. \label{eq:projection_matrix}
\end{equation}
Define the commanded non-gravitational transverse acceleration by
\begin{equation}
\mathbf a^{\perp}_{\mathrm{cmd},E}
=
P_{\perp,E}\bigl(\mathbf a_{\mathrm{cmd},E}-\mathbf g_E\bigr).
\end{equation}
Resolve this command into the two transverse body axes:
\begin{equation}
\mathbf a^{\perp}_{\mathrm{cmd},B}
=
\begin{bmatrix}
\hat{\mathbf e}_{1,B}^{\mathsf T}\\[2pt]
\hat{\mathbf e}_{2,B}^{\mathsf T}
\end{bmatrix}
C_{BE}\mathbf a^{\perp}_{\mathrm{cmd},E}
=
\begin{bmatrix}
a_{\mathrm{cmd},1}\\[2pt]
a_{\mathrm{cmd},2}
\end{bmatrix}. \label{eq:accel_body_resolved}
\end{equation}
The allocator is then
\begin{equation}
\begin{bmatrix}
\delta_{1,\mathrm{cmd}}\\[2pt]
\delta_{2,\mathrm{cmd}}
\end{bmatrix}
=
\mathcal A\!\left(
a_{\mathrm{cmd},1},
a_{\mathrm{cmd},2},
\alpha,
M_\infty,
q_\infty
\right). \label{eq:allocator_func}
\end{equation}
Here \(a_{\mathrm{cmd},1}\) and \(a_{\mathrm{cmd},2}\) are simply the desired non-gravitational acceleration components along the body axes \(+\hat{\mathbf e}_{1,B}\) and \(+\hat{\mathbf e}_{2,B}\). For example, \(a_{\mathrm{cmd},1}>0\) asks the control system to generate acceleration in the \(+\hat{\mathbf e}_{1,B}\) direction. The allocator \(\mathcal A\) should therefore be documented with the same sign convention: a positive output deflection should mean the right-hand-rule rotation about the corresponding hinge axis introduced in Step~4. (\textbf{Implementation Note:} strictly verify this matches the mathematical signs below; sign-flip bugs between allocator and aerodynamics are a frequent source of closed-loop instability!)

Thus, the attitude information needed by the allocator is not passed as a separate tracking variable; it is naturally contained in the transformed command \(\mathbf a^{\perp}_{\mathrm{cmd},B}\) via the current attitude matrix \(C_{BE}\).

\subsection{Step 3: Transform the actual air-relative velocity into body coordinates}

The actual air-relative velocity used by the aerodynamic model is
\begin{equation}
\mathbf V_{\mathrm{rel},E}
=
\mathbf V_{\mathrm{wind},E}-\mathbf v_E,
\end{equation}
and its body-frame components are
\begin{equation}
\mathbf V_{\mathrm{rel},B}
=
C_{BE}\mathbf V_{\mathrm{rel},E},
\qquad
\hat{\mathbf u}_B
=
\frac{\mathbf V_{\mathrm{rel},B}}
{\|\mathbf V_{\mathrm{rel},B}\|}. \label{eq:rel_wind_step3}
\end{equation}
This vector already contains the combined effects of actual angle of attack, geometric banking, and wind. As a sign check, if the atmosphere is still then \(\mathbf V_{\mathrm{rel},E}=-\mathbf v_E\); a tailwind parallel to the vehicle motion makes \(\|\mathbf V_{\mathrm{rel},E}\|\) smaller, whereas a headwind makes it larger. \textbf{Implementation Note:} Here again is a vector normalization that may require care in the low-speed regime.

\subsection{Step 4: Compute the incidence of the wind on each physical flap}

The pair commands \((\delta_1,\delta_2)\) remain the actuator inputs, but the aerodynamic loading should be evaluated flap-by-flap. Let \(i\in\{1,2,3,4\}\) index the four physical flaps, each of area \(A_f\). \textbf{Implementation Note:} Ensure array indexing in your software maps correctly to these 1-based formulas. 

For a cone with half-angle \(\theta_c\), the inward-pointing undeflected flap normals may be written as
\begin{equation}
\begin{aligned}
\hat{\mathbf n}_{10,B}
&=
-\cos\theta_c\,\hat{\mathbf e}_{1,B}
+\sin\theta_c\,\hat{\mathbf e}_{3,B}, \\
\hat{\mathbf n}_{20,B}
&=
-\cos\theta_c\,\hat{\mathbf e}_{2,B}
+\sin\theta_c\,\hat{\mathbf e}_{3,B}, \\
\hat{\mathbf n}_{30,B}
&=
+\cos\theta_c\,\hat{\mathbf e}_{1,B}
+\sin\theta_c\,\hat{\mathbf e}_{3,B}, \\
\hat{\mathbf n}_{40,B}
&=
+\cos\theta_c\,\hat{\mathbf e}_{2,B}
+\sin\theta_c\,\hat{\mathbf e}_{3,B}.
\end{aligned} \label{eq:flap_normals}
\end{equation}
If the actual flap mounting differs from this idealized cone geometry, these four normals should be replaced by geometry-derived or CAD-derived values.

Let the pair-1 flaps share the hinge axis
\begin{equation}
\hat{\mathbf h}_{1,B}=\hat{\mathbf h}_{3,B}=\hat{\mathbf e}_{2,B},
\end{equation}
and let the pair-2 flaps share the hinge axis
\begin{equation}
\hat{\mathbf h}_{2,B}=\hat{\mathbf h}_{4,B}=-\hat{\mathbf e}_{1,B}.
\end{equation}
The commanded pair deflections map to the four physical flaps as
\begin{equation}
\delta_{f,1}=\delta_{f,3}=\delta_1,
\qquad
\delta_{f,2}=\delta_{f,4}=\delta_2.
\end{equation}
The deflected normal of flap \(i\) is then
\begin{equation}
\hat{\mathbf n}_{i,B}(\delta_{f,i})
=
R(\hat{\mathbf h}_{i,B},\delta_{f,i})\hat{\mathbf n}_{i0,B},
\qquad
i\in\{1,2,3,4\}, \label{eq:flap_deflected_normal}
\end{equation}
where \(R(\hat{\mathbf h},\delta)\) is the Rodrigues rotation operator,
\begin{equation}
R(\hat{\mathbf h},\delta)\mathbf a
=
\mathbf a\cos\delta
+
(\hat{\mathbf h}\times\mathbf a)\sin\delta
+
\hat{\mathbf h}\left(\hat{\mathbf h}\cdot\mathbf a\right)(1-\cos\delta). \label{eq:rodrigues_rotation}
\end{equation}

This is just the standard three-dimensional right-hand-rule rotation of the vector \(\mathbf a\) through the angle \(\delta\) about the axis \(\hat{\mathbf h}\). For example, because \(\hat{\mathbf h}_{1,B}=+\hat{\mathbf e}_{2,B}\), a positive \(\delta_1\) means a right-hand-rule rotation about \(+\hat{\mathbf e}_{2,B}\) for both flaps 1 and 3. (\textbf{Implementation Note:} You may not need to implement Rodrigues' formula from scratch. Many 3D spatial libraries have built-in functions to generate a rotation from an ``axis and angle''. You can simply pass the hinge axis $\hat{\mathbf h}_{i,B}$ and deflection angle $\delta_{f,i}$ to your library's rotation object and apply it directly to the undeflected normal $\hat{\mathbf n}_{i0,B}$. However, when using external rotation libraries, strictly verify their mapping convention (active vs passive) to ensure it acts correctly on the vector.)

Define the loaded-side incidence factor of each flap by
\begin{equation}
\lambda_i(\delta_{f,i})
=
\max\!\left(0,\hat{\mathbf u}_B\cdot \hat{\mathbf n}_{i,B}(\delta_{f,i})\right),
\qquad
i\in\{1,2,3,4\}. \label{eq:incidence_factor}
\end{equation}
(Note: The \(\max()\) operator is \(C^0\) continuous but not differentiable at \(0\). If this model is ever to be used in trajectory optimization or Unscented Kalman Filter tracking schemes requiring gradient availability, a smooth soft-max substitution might be required.)
The corresponding local incidence angles are
\begin{equation}
\beta_i=\sin^{-1}\lambda_i(\delta_{f,i}),
\qquad
\beta_{i0}=\sin^{-1}\lambda_i(0). \label{eq:incidence_angle}
\end{equation}
\textbf{Implementation Note:} In code, clamp \(\lambda_i\) to the interval \([0,1]\) before evaluating \(\sin^{-1}\), since floating-point round-off can otherwise produce \(\lambda_i>1\).

The \(\max(0,\cdot)\) makes the model one-sided: only the windward side of a flap carries pressure in this simplified treatment. Thus \(\lambda_i>0\) means the inward-pointing normal has a positive component toward the oncoming air and the flap is loaded on that side; \(\lambda_i=0\) means the chosen pressure side is turned away from the wind. Because \(\hat{\mathbf u}_B\) already contains the full wind direction in the body frame, these expressions automatically account for the current angle of attack, the local wind azimuth, and the commanded pair deflections. My original lumped-pair construction may still be used as a lower-fidelity approximation, but it does not distinguish the different loading of the two opposite flaps within a pair, which will matter if the angle of attack becomes large.

\subsection{Step 5: Compute the incremental force vector of each flap}

Van Brunt tabulates the undeflected-body coefficients \(C_D(12)\), \(C_L(12)\), \(C_N(12)\), \(C_M(12)\), and \(C_{M_q}(12)\), but he does \emph{not} separately tabulate \(C_{p,\max}(12)\). Accordingly, the flap model should not assume that such a tabulated value exists. Instead, the flap-force model needs its own pressure-scale factor. A simple modified-Newtonian choice is
\begin{equation}
K_f(M_\infty)=C_{p,\max}(M_\infty), \label{eq:kf_pressure_scale}
\end{equation}
where \(C_{p,\max}(M_\infty)\) is computed from the stagnation-pressure relation introduced earlier in Section~3. 

With that independently chosen pressure scale, the absolute force on physical flap \(i\) is approximated by
\begin{equation}
\mathbf F_{f,i,B}^{(\mathrm{abs})}(\delta_{f,i})
=
q_\infty A_f K_f(M_\infty)\,
\lambda_i(\delta_{f,i})^2\,
\hat{\mathbf n}_{i,B}(\delta_{f,i}),
\qquad
i\in\{1,2,3,4\}. \label{eq:flap_abs_force}
\end{equation}
Because \(K_f>0\) and \(\lambda_i^2\ge 0\), the sign of the vector force comes entirely from the chosen normal direction \(\hat{\mathbf n}_{i,B}\). For example, a loaded flap 1 produces a force with the same general sign pattern as \(\hat{\mathbf n}_{10,B}=-\cos\theta_c\,\hat{\mathbf e}_{1,B}+\sin\theta_c\,\hat{\mathbf e}_{3,B}\): inward toward the axis and aft toward the base.

Because the undeflected flaps are already included in the baseline aerodynamic tables, only the \emph{incremental} flap force should be added:
\begin{equation}
\Delta\mathbf F_{f,i,B}
=
\mathbf F_{f,i,B}^{(\mathrm{abs})}(\delta_{f,i})
-
\mathbf F_{f,i,B}^{(\mathrm{abs})}(0). \label{eq:flap_inc_force}
\end{equation}
This subtraction permits both additional loading and unloading relative to the baseline state.

The incremental drag and lift vectors of flap \(i\) are obtained by decomposition along and normal to the freestream:
\begin{equation}
\Delta\mathbf D_{f,i,B}
=
\left(
\Delta\mathbf F_{f,i,B}\cdot \hat{\mathbf u}_B
\right)\hat{\mathbf u}_B,
\qquad
\Delta\mathbf L_{f,i,B}
=
\Delta\mathbf F_{f,i,B}
-
\Delta\mathbf D_{f,i,B}.
\end{equation}
Here \(+\hat{\mathbf u}_B\) is the oncoming-air direction, so a positive drag component acts opposite the vehicle motion in still air. If scalar magnitudes are desired, they reduce to
\begin{align}
\Delta D_{f,i}
&=
q_\infty A_f K_f(M_\infty)
\left[
\sin^3\beta_i-\sin^3\beta_{i0}
\right], \\
\Delta L_{f,i}
&=
q_\infty A_f K_f(M_\infty)
\left[
\sin^2\beta_i\cos\beta_i
-
\sin^2\beta_{i0}\cos\beta_{i0}
\right].
\end{align}
(These scalar formulas are provided for intuition, but directly computing the vector subtraction above is usually cleaner in code.)
The total incremental flap force is then
\begin{equation}
\Delta\mathbf F_{f,B}
=
\sum_{i=1}^4 \Delta\mathbf F_{f,i,B}. \label{eq:total_flap_force}
\end{equation}

\subsection{Step 6: Compute the baseline body force from the aerodynamic tables}

Interpolate the undeflected-vehicle tables at the current \((\alpha,M_\infty)\), then evaluate
\begin{equation}
\mathbf F_{\mathrm{body},B}
=
q_\infty S_{\mathrm{ref}}
\left[
C_D(\alpha,M_\infty)\hat{\mathbf u}_B
+
C_L(\alpha,M_\infty)\hat{\boldsymbol\ell}_B
\right]. \label{eq:step6_body_force}
\end{equation}
The sign convention is again carried by the basis vectors. A positive \(C_D\) produces force along \(+\hat{\mathbf u}_B\), and a positive \(C_L\) produces force along \(+\hat{\boldsymbol\ell}_B\); if a coefficient is negative, the force automatically reverses direction. If desired, \(C_N(\alpha,M_\infty)\) may be used as a cross-check because the body is nearly axisymmetric, but only one force representation should be inserted into the equations of motion.

\subsection{Step 7: Sum the aerodynamic forces and return them to ECI}

The total aerodynamic force in body coordinates is
\begin{equation}
\mathbf F_{\mathrm{aero},B}
=
\mathbf F_{\mathrm{body},B}
+
\Delta\mathbf F_{f,B}
=
\mathbf F_{\mathrm{body},B}
+
\sum_{i=1}^4 \Delta\mathbf F_{f,i,B}. \label{eq:total_aero_force_b}
\end{equation}
Transform this back to ECI using the constructed mapping matrix:
\begin{equation}
\mathbf F_{\mathrm{aero},E}
=
C_{EB}\mathbf F_{\mathrm{aero},B}. \label{eq:total_aero_force_e}
\end{equation}
The translational update is then
\begin{equation}
m\dot{\mathbf v}_E
=
m\mathbf g_E(\mathbf r_E)+\mathbf F_{\mathrm{aero},E},
\qquad
\dot{\mathbf r}_E=\mathbf v_E.
\end{equation}

\subsection{Step 8: Compute the total aerodynamic moment in body coordinates}

The baseline body moment is
\begin{equation}
\mathbf M_{\mathrm{body},B}
=
q_\infty S_{\mathrm{ref}}c_{\mathrm{ref}}
\,C_M(\alpha,M_\infty)\hat{\mathbf m}_B
+
q_\infty S_{\mathrm{ref}}c_{\mathrm{ref}}
\left(\frac{c_{\mathrm{ref}}}{2V}\right)
C_{M_q}(\alpha,M_\infty)\,\boldsymbol\omega_{\perp,B}. \label{eq:step8_body_moment}
\end{equation}
This uses the tabulated \(C_M\) directly, which is preferable to a linearized \(C_{M_\alpha}\alpha\) approximation when the tables are available.

Treat each physical flap as a point force applied at its own centroid. Let \(r_{f,\perp}>0\) denote the radial distance from the symmetry axis to the flap centroid in the transverse plane, and let \(c_f>0\) denote the aft distance from the center of mass to the flap centroid measured in the \(+\hat{\mathbf e}_{3,B}\) direction. With the present convention, \(+\hat{\mathbf e}_{3,B}\) points in the zero-angle-of-attack oncoming-flow direction and therefore toward the base of the vehicle. Then a convenient idealization is
\begin{equation}
\begin{aligned}
\mathbf r_{f,1,B}
&=
+r_{f,\perp}\,\hat{\mathbf e}_{1,B}
+c_f\,\hat{\mathbf e}_{3,B}, \\
\mathbf r_{f,2,B}
&=
+r_{f,\perp}\,\hat{\mathbf e}_{2,B}
+c_f\,\hat{\mathbf e}_{3,B}, \\
\mathbf r_{f,3,B}
&=
-r_{f,\perp}\,\hat{\mathbf e}_{1,B}
+c_f\,\hat{\mathbf e}_{3,B}, \\
\mathbf r_{f,4,B}
&=
-r_{f,\perp}\,\hat{\mathbf e}_{2,B}
+c_f\,\hat{\mathbf e}_{3,B}.
\end{aligned} \label{eq:flap_centroids}
\end{equation}
If a different body-axis convention is adopted in code, these signs should be changed consistently with that convention.
The incremental aerodynamic moment of flap \(i\) is
\begin{equation}
\Delta\mathbf M_{f,i,B}
=
\mathbf r_{f,i,B}\times \Delta\mathbf F_{f,i,B},
\qquad
i\in\{1,2,3,4\}. \label{eq:flap_moment}
\end{equation}
Thus,
\begin{equation}
\mathbf M_{\mathrm{aero},B}
=
\mathbf M_{\mathrm{body},B}
+
\sum_{i=1}^4 \Delta\mathbf M_{f,i,B}. \label{eq:total_aero_moment}
\end{equation}
This flap-wise form preserves the net pitch and yaw moments under asymmetric flap loading. In a full three-axis rotational model it also produces the corresponding axial-roll moment. In the present reduced model, however, that axial component is intentionally discarded by the constraint \(\omega_{3,B}=0\), \(M_{3,B}=0\).\footnote{If one intentionally wants to move back to a control vector and Euler method with orthonormalization instead of using quaternions, the roll moment is recovered by setting \(r_{f,\perp}=0\) and combining opposite flaps after their individual forces are summed.}

If flap loading must be capped explicitly, it is cleaner to clip the force on each physical flap before computing the moment:
\begin{equation}
\Delta\mathbf F_{f,i,B}
\leftarrow
\operatorname{clip}
\!\left(
\|\Delta\mathbf F_{f,i,B}\|,
0,
F_{\mathrm{flap,max}}
\right)
\frac{\Delta\mathbf F_{f,i,B}}
{\|\Delta\mathbf F_{f,i,B}\|}, \label{eq:flap_clipping}
\end{equation}
with the obvious convention that a zero vector remains zero. (\textbf{Implementation Note:} strictly handle the condition where \(\|\Delta\mathbf F_{f,i,B}\| < \epsilon\) prior to normalizing to avoid division-by-zero or \texttt{NaN} state pollution!) The moment is then recomputed from \(\Delta\mathbf M_{f,i,B}=\mathbf r_{f,i,B}\times \Delta\mathbf F_{f,i,B}\). If a moment cap rather than a force cap is available, use \(F_{\mathrm{flap,max}}=M_{\mathrm{flap,max}}/\|\mathbf r_{f,i,B}\|\) as the local equivalent limit.

\subsection{Step 9: Update the rotational state over the time step}
The full Euler equations are

\begin{align}
M_{1,B} &= I_\perp \dot{\omega}_{1,B} + (I_3-I_\perp)\omega_{2,B}\omega_{3,B},\\
M_{2,B} &= I_\perp \dot{\omega}_{2,B} - (I_3-I_\perp)\omega_{1,B}\omega_{3,B},\\
M_{3,B} &= I_3 \dot{\omega}_{3,B}.
\end{align}
However, because commanded and induced body-axial twist are to be neglected in flight dynamics, we may enforce \(\omega_{3,B}=0\) to produce the reduced model
\begin{align}
M_{1,B} &= I_\perp \dot{\omega}_{1,B},\\
M_{2,B} &= I_\perp \dot{\omega}_{2,B},\\
M_{3,B} &= 0.
\end{align}

Here \(M_{k,B}\) and \(\omega_{k,B}\) are simply the components of aerodynamic moment and angular velocity about the body axes \(\hat{\mathbf e}_{k,B}\). 

Propagate the unit quaternion \(\mathbf q_{EB}^n\), where the superscript \(n\) represents the state necessary for the \(n^\textrm{th}\) time step. Over a single step, assume \(\boldsymbol\omega_B\) is frozen at its step value \(\boldsymbol\omega_B^n\), define
\begin{equation}
\theta^n = \|\boldsymbol\omega_B^n\|\,\Delta t,
\end{equation}
and form the body-frame incremental quaternion
\begin{equation}
\Delta\mathbf q_B^n
=
\begin{bmatrix}
\cos(\theta^n/2)\\[2pt]
\dfrac{\boldsymbol\omega_B^n}{\|\boldsymbol\omega_B^n\|}\sin(\theta^n/2)
\end{bmatrix}, \label{eq:delta_quat}
\end{equation}
recalling that placing the three spatial components of \(\boldsymbol{\omega}_B^n\) in the lower block expands the expression to the full \(4\times1\) quaternion vector. \textbf{Implementation Note:} To avoid division by \(\|\boldsymbol\omega_B^n\|\) when the step angle is small, it is numerically cleaner to evaluate
\begin{equation}
\Delta\mathbf q_B^n
=
\begin{bmatrix}
\cos(\theta^n/2)\\[2pt]
\frac{\Delta t}{2}\,\operatorname{sinc}\!\left(\frac{\theta^n}{2}\right)\boldsymbol\omega_B^n
\end{bmatrix},
\qquad
\operatorname{sinc}(x)\equiv\frac{\sin x}{x},
\end{equation}
with the continuous value \(\operatorname{sinc}(0)=1\).\footnote{These are equivalent expressions:
\begin{align*}
\Delta\mathbf q_B^n
&=
\begin{bmatrix}
\cos(\theta^n/2)\\[2pt]
\dfrac{\boldsymbol\omega_B^n}{\|\boldsymbol\omega_B^n\|}\sin(\theta^n/2)
\end{bmatrix},
\qquad
\theta^n = \|\boldsymbol\omega_B^n\|\Delta t
\\[6pt]
&=
\begin{bmatrix}
\cos(\theta^n/2)\\[2pt]
\dfrac{\boldsymbol\omega_B^n}{\|\boldsymbol\omega_B^n\|}\,
\operatorname{sinc}\!\left(\frac{\theta^n}{2}\right)\,
\frac{\theta^n}{2}
\end{bmatrix}
\qquad
\left(
\operatorname{sinc}(x)\equiv \frac{\sin x}{x}
\right)
\\[6pt]
&=
\begin{bmatrix}
\cos(\theta^n/2)\\[2pt]
\dfrac{\boldsymbol\omega_B^n}{\|\boldsymbol\omega_B^n\|}\,
\operatorname{sinc}\!\left(\frac{\theta^n}{2}\right)\,
\frac{\|\boldsymbol\omega_B^n\|\Delta t}{2}
\end{bmatrix}
\\[6pt]
&=
\begin{bmatrix}
\cos(\theta^n/2)\\[2pt]
\frac{\Delta t}{2}\,
\operatorname{sinc}\!\left(\frac{\theta^n}{2}\right)\,
\boldsymbol\omega_B^n
\end{bmatrix}.
\end{align*}
}


%If \(\|\boldsymbol\omega_B^n\|\) is numerically very small (e.g., using a tolerance bound \(\|\boldsymbol\omega_B^n\| < 10^{-6}\) in code), replace the trigonometric form by its small-angle limit:
%\begin{equation}
%\Delta\mathbf q_B^n
%\approx
%\begin{bmatrix}
%1\\[2pt]
%\tfrac12\boldsymbol\omega_B^n\Delta t
%\end{bmatrix}, \label{eq:small_angle_quat}
%\end{equation}
%which avoids division by a tiny number. 

Using the standard Hamilton quaternion product\footnote{\url{https://www.mathworks.com/help/aeroblks/quaternionmultiplication.html}} \(\otimes\), the discrete attitude update is then
\begin{equation}
\mathbf q_{EB}^{n+1}
=
\mathbf q_{EB}^{n}\otimes \Delta\mathbf q_B^n. \label{eq:quat_update}
\end{equation}
After the multiplication, renormalize the quaternion:
\begin{equation}
\mathbf q_{EB}^{n+1}
\leftarrow
\frac{\mathbf q_{EB}^{n+1}}{\|\mathbf q_{EB}^{n+1}\|}.
\end{equation}
The rotation matrix $C_{EB}$ and fundamental body axis vector can then be reconstructed as shown in equation~\ref{eq:quat_to_dcm} when queried by the simulation.

This quaternion step is the exact constant-rate update over one time step and is the preferred implementation because the only geometric correction required is the cheap normalization of \(\mathbf q_{EB}\).


\subsection{Step 10: Recompute the actual angle of attack from the updated state}

After the body-frame has been evolved, we need to recompute the angle of attack. Using the recovered \(C_{EB}^{n+1}\) from Equation \eqref{eq:quat_to_dcm}, set \(C_{BE}^{n+1}=\left(C_{EB}^{n+1}\right)^{\mathsf T}\). Then
\begin{equation}
\mathbf V_{\mathrm{rel},E}^{\,n+1}
=
\mathbf V_{\mathrm{wind},E}^{\,n+1}
-
\mathbf v_E^{\,n+1},
\qquad
\hat{\mathbf u}_B^{\,n+1}
=
C_{BE}^{\,n+1}
\frac{\mathbf V_{\mathrm{rel},E}^{\,n+1}}
{\|\mathbf V_{\mathrm{rel},E}^{\,n+1}\|}.
\end{equation}
Then
\begin{equation}
\alpha^{n+1}
=
\cos^{-1}\!\left(u_3^{n+1}\right),
\qquad
\chi_\alpha^{n+1}
=
\operatorname{atan2}\!\left(u_2^{n+1},u_1^{n+1}\right),
\end{equation}
where \(\alpha\) is the unsigned magnitude of the total angle of attack and \(\chi_\alpha\) carries the azimuthal sign information in the body-transverse plane. 

\textbf{Implementation Notes:} The azimuth \(\chi_\alpha\) is undefined when \(u_1=u_2=0\). In code, if \( s=\sqrt{u_1^2+u_2^2}<\epsilon_s,\) set \(\chi_\alpha\) to zero or retain its previous value only as a placeholder for continuity. Additionally, to prevent numerical errors from round-offs, clamp \(u_3\) before taking the inverse cosine \(u_3\leftarrow\textrm{clip}(u_3,-1,1)\).

The corresponding two-component attack vector is
\begin{equation}
\boldsymbol\alpha_{\perp,B}^{\,n+1}
=
\frac{1}{u_3^{n+1}}
\begin{bmatrix}
u_1^{n+1}\\[2pt]u_2^{n+1}
\end{bmatrix}.
\end{equation}
This quantity is the correct ``vector sum of pitch and yaw'' for the reduced model.
\textbf{Implementation Note:} This vector is useful near trim, but it should only be formed when \(|u_3|>\epsilon_u\). If \(|u_3|\le \epsilon_u\), the vehicle is near a \(90^\circ\) angle-of-attack or tumbling state outside the intended regime of the reduced model; in that case use \((\alpha,\chi_\alpha)\) or the full unit-wind vector \(\hat{\mathbf u}_B\) directly rather than \(\boldsymbol\alpha_{\perp,B}\).

\subsection{Summary}
The resulting loop is therefore:
\begin{center}
\texttt{guidance command in ECI}
\(\to\)
\texttt{body-frame allocator}
\(\to\)
\texttt{pair loads in \(\mathcal B\)}
\\[2pt]
\texttt{body force and moment sums}
\(\to\)
\texttt{ECI translation + quaternion attitude update}
\(\to\)
\texttt{recomputed \((\alpha,\chi_\alpha)\)}.
\end{center}
This closes the simulation natively tracking dynamic 3D geometric attitude without requiring a full 6-DOF mapping from commanded acceleration to explicit rigid-body roll variables.

\bibliographystyle{plain}
\bibliography{refs}
\clearpage
\appendix
\section{Estimation of Real-Gas Corrections}\label{app:realgas}

To extend the Mach-12 MC-NEW baseline of Van Brunt to other hypersonic conditions, we define the Mach correction multiplier
\begin{equation}
r(M_\infty,h)\equiv \frac{C_{p,\max}(M_\infty,h)}{C_{p,\max}(12,h)}, \label{eq:rg_mach_multiplier}
\end{equation}

To improve the accuracy, a real-gas workup was performed. The maximum pressure coefficient (and thus its dependent coefficients such as \(C_D\) and \(C_L\)), can then be evaluated as
\begin{equation}
C_{p,\max}(M_\infty,h)=\frac{p_{0,2}(M_\infty,h)-p_\infty(h)}{q_\infty(h)},
\qquad
q_\infty(h)=\frac{1}{2}\rho_\infty(h)V_\infty^2, \label{eq:rg_cp_max}
\end{equation}
rather than from a constant-\(\gamma\) perfect-gas formula.

For each tuple \((M_\infty,h)\), the upstream freestream state \((p_\infty,T_\infty,\rho_\infty)\) was taken from a standard-atmosphere model \cite{nasa_atmosphere}, the freestream speed was set by \(V_\infty=M_\infty a_\infty\), and the dynamic pressure was formed as
\begin{equation}
q_\infty=\frac{1}{2}\rho_\infty V_\infty^2.
\end{equation} 
The shock calculation itself followed the normal-shock/stagnation-pressure framework of NACA Report 1135 \cite{naca1135}, but the perfect-gas closure was replaced by an equilibrium-air thermodynamic closure using NASA Glenn thermodynamic polynomial data \cite{mcbride2002} for the species $\textrm{N}_2$, $\textrm{O}_2$, $\textrm{NO}$, O, N, $\textrm{N}_2^+$, $\textrm{O}_2^+$, $\textrm{NO}^+$, $\textrm{N}^+$, $\textrm{O}^+$, and $e^-$. In practice, this means that the mixture thermodynamics were evaluated with temperature-dependent species properties and equilibrium composition, so the calculation includes variable \(c_p\), variable effective \(\gamma\), vibrational excitation, equilibrium dissociation, and equilibrium ionization, in the same general sense as a CEA-type equilibrium calculation \cite{gordon1994cea,leader2024cea}. 

It is important to note what was \emph{not} included. The calculation treated the gas as being in chemical equilibrium and therefore did not model finite-rate nonequilibrium chemistry, viscous or boundary-layer interaction, catalytic wall effects, or radiation. Accordingly, the tabulated \(r(M_\infty,h)\) values should be interpreted as an equilibrium-thermochemistry correction to the pressure scale, not as a full high-enthalpy CFD solution \cite{gordon1994cea,leader2024cea}. 

The resulting values of \(r(M_\infty,h)\) are shown below in Table~\ref{tab:realgascorrection} for integer Mach numbers from 5 to 25 and for three representative altitude strata: sea level, \(40{,}000\) ft, and \(80{,}000\) ft.

\begin{table}[h!]
\centering
\caption{Real-gas Mach correction multiplier \(r(M_\infty,h)\) normalized to the Mach-12 baseline.}\label{tab:realgascorrection}
\begin{tabular}{c|ccc}
\hline
Mach & \(r\) at 0 kft & \(r\) at 40 kft & \(r\) at 80 kft \\
\hline
5  & 0.9658 & 0.9642 & 0.9612 \\
6  & 0.9739 & 0.9720 & 0.9690 \\
7  & 0.9802 & 0.9779 & 0.9750 \\
8  & 0.9856 & 0.9829 & 0.9804 \\
9  & 0.9905 & 0.9876 & 0.9861 \\
10 & 0.9947 & 0.9923 & 0.9917 \\
11 & 0.9979 & 0.9966 & 0.9964 \\
12 & 1.0000 & 1.0000 & 1.0000 \\
13 & 1.0013 & 1.0025 & 1.0024 \\
14 & 1.0023 & 1.0041 & 1.0036 \\
15 & 1.0034 & 1.0050 & 1.0040 \\
16 & 1.0048 & 1.0057 & 1.0048 \\
17 & 1.0063 & 1.0067 & 1.0062 \\
18 & 1.0078 & 1.0080 & 1.0078 \\
19 & 1.0091 & 1.0094 & 1.0094 \\
20 & 1.0104 & 1.0108 & 1.0110 \\
21 & 1.0114 & 1.0122 & 1.0124 \\
22 & 1.0123 & 1.0134 & 1.0136 \\
23 & 1.0131 & 1.0145 & 1.0147 \\
24 & 1.0136 & 1.0154 & 1.0157 \\
25 & 1.0141 & 1.0163 & 1.0165 \\
\hline
\end{tabular}
\end{table}

The corresponding equilibrium-air reference values used for normalization are
\[
C_{p,\max}(12)=
\begin{cases}
1.8927, & 0\ \text{kft},\\
1.8897, & 40\ \text{kft},\\
1.8960, & 80\ \text{kft}.
\end{cases}
\]

For implementation in a control or trajectory code, one may use
\begin{equation}
C_p(M_\infty,h)\approx r(M_\infty,h)\,C_p(12,h),
\end{equation}
with bilinear interpolation in Mach and altitude. Because the dependence on altitude is weak, an altitude-independent approximation may be preferred, using the \(80\)-kft column as a reasonable single-altitude schedule.
\end{document}